\documentclass[11pt]{article}

\usepackage[a4paper,margin=1in]{geometry}
\usepackage[T1]{fontenc}
\usepackage{lmodern}
\usepackage{microtype}
\usepackage{amsmath,amssymb,amsthm,mathtools}
\usepackage{booktabs,tabularx}
\usepackage{enumitem}
\usepackage[numbers,sort&compress]{natbib}
\usepackage[colorlinks=true,linkcolor=blue,citecolor=blue,urlcolor=blue]{hyperref}
\usepackage[nameinlink,noabbrev]{cleveref}

\setlength{\emergencystretch}{3em}

\newtheorem{theorem}{Theorem}[section]
\newtheorem{proposition}[theorem]{Proposition}
\newtheorem{corollary}[theorem]{Corollary}
\newtheorem{lemma}[theorem]{Lemma}
\theoremstyle{definition}
\newtheorem{definition}[theorem]{Definition}
\theoremstyle{remark}
\newtheorem{remark}[theorem]{Remark}

\newcommand{\one}{\mathbf 1}
\newcommand{\Lzero}{\mathrm{L0}}
\newcommand{\Lone}{\mathrm{L1}}
\newcommand{\Ltwo}{\mathrm{L2}}

\title{\textbf{Exceptional-Scale Selection and the Actual
Fixed-Two Power Gate}\\
\large A Carleson Return Interface for H3}
\author{Liang Wang}
\date{July 2026}

\begin{document}
\maketitle

\begin{abstract}
An arithmetic theorem valid at almost all logarithmic scales need
not control a deterministically selected physical prefix.  We give
two noncircular return interfaces for a fixed-two TPC packet.  The
first is a direct uniform all-prefix theorem.  The second requires
the actual selector measure \(\nu_X\) to be dominated by normalized
logarithmic measure \(m_X\):
\[
 \nu_X(E)\le K_Xm_X(E)
\]
for every Borel set \(E\).  An arithmetic scale-average error
\(\varepsilon_X\) then returns as \(K_X\varepsilon_X\).  Combining
this with squarefree-tail, component-census and bounded-variation
losses gives the exact amplitude ledger
\[
 \sigma_{\rm raw}
 =
 \min\{\eta_{\rm tail},
 \sigma_{\rm aff}-\ell_{\rm sel}-\ell_{\rm cen}-\ell_{\rm BV}\}.
\]
TPC-131 closes only if
\(\sigma_{\rm raw}>\Lambda_{\rm phys}\).  Atomic selectors are
singular with respect to continuous logarithmic measure, so a
discrete all-prefix family requires a pointwise theorem or a proved
smoothing crosswalk.  Current fixed-data logarithmic inputs have no
positive \(X\)-power exponent.  Tao--Ter\"av\"ainen's 2026
small-polylog affine theorem does add a genuine power-of-log
almost-scale corridor.  Its cumulative exceptional-set estimate is
not automatically the same power-of-log estimate on an arbitrary
terminal window: that conversion has a separate exact loss.  After
this window loss, selector loss and reassembly losses, it can yield
qualitative cancellation only if the full logarithmic exponent
ledger stays positive.  It still has \(X\)-power exponent zero.
The actual selector and all-prefix gates remain open; no positive
\(\Ltwo\), H3, or \(1/400\) conclusion is claimed.
\end{abstract}

\section{Why almost all scales are not all prefixes}

For a bounded arithmetic coefficient \(b(n)\), write
\[
 C(t)=\frac1t\sum_{t<n\le2t}b(n).
\]
An almost-scale theorem has the shape
\begin{equation}
 \frac1{\log\omega}
 \int_{x/\omega}^{x}|C(t)|\frac{dt}{t}
 \le\varepsilon_X.
\label{eq:almost-scale}
\end{equation}
Helfgott and Radziwi{\l}{\l} prove a quantitative almost-all-scale
conclusion for the consecutive Liouville correlation
\citep{HelfgottRadziwill2021}.  We use \eqref{eq:almost-scale} as an
abstract normalized interface rather than attributing this exact
continuous formulation to that source.  TPC-138 separates its
even-carrier shift-one consequence from the active odd carrier and
also records the newer Tao--Ter\"av\"ainen small-polylog affine
corridor \citep{WangTPC138,TaoTeravainen2026}.  The remaining
logical issue is independent of determinant: a theorem outside a
small exceptional set is still not a deterministic all-prefix
theorem.  Moreover, a cumulative exceptional-set bound up to \(x\)
is not automatically a normalized bound on every shorter terminal
window.

TPC-130 requires the declared actual block-prefix family
\citep{WangTPC130}.  A deterministic list of its endpoints may lie
inside the exceptional scale set unless a return theorem prevents
that concentration.

\begin{proposition}[Average-to-maximum inference is invalid]
\label{prop:no-max}
From \eqref{eq:almost-scale} alone one cannot deduce
\[
 \sup_{x/\omega\le t\le x}|C(t)|=o(1).
\]
The failure persists even for bounded measurable \(C\).
\end{proposition}

\begin{proof}
Choose a set \(E_X\) of normalized logarithmic measure
\(\varepsilon_X\), and put \(C=\one_{E_X}\).  Then
\eqref{eq:almost-scale} holds, while the supremum is one.
\end{proof}

This counterexample is analytic, not a construction of Liouville
values.  It proves that an additional selector or pointwise premise
is logically necessary.

\section{Two legal return interfaces}

Let
\[
 J_X=[x/\omega,x],\qquad
 m_X(E)=\frac1{\log\omega}\int_{E\cap J_X}\frac{dt}{t}.
\]
Thus \(m_X\) is a probability measure on \(J_X\).

\begin{proposition}[Global exceptional density versus a terminal window]
\label{prop:global-to-window}
Suppose an exceptional set \(\mathcal E\subset[1,\infty)\) satisfies
\begin{equation}
 \frac1{\log y}
 \int_{[1,y]\cap\mathcal E}\frac{dt}{t}
 \ll(\log y)^{-c_0}
 \qquad(y\ge2).
\label{eq:global-exceptional}
\end{equation}
Then
\begin{equation}
 \boxed{
 m_X(\mathcal E)
 \ll
 \min\left\{
 1,\frac{(\log x)^{1-c_0}}{\log\omega}
 \right\}.}
\label{eq:window-exceptional}
\end{equation}
In particular, \eqref{eq:global-exceptional} alone does not give
\(m_X(\mathcal E)\ll(\log x)^{-c_0}\) for an arbitrary
\(\omega=\omega(x)\).
\end{proposition}

\begin{proof}
The numerator defining \(m_X(\mathcal E)\) is bounded by the
integral over \([1,x]\cap\mathcal E\), which is
\(O((\log x)^{1-c_0})\) by
\eqref{eq:global-exceptional}.  Divide by \(\log\omega\), and also
use \(m_X(\mathcal E)\le1\).
\end{proof}

\begin{remark}[Window exponent]
\label{rem:window-exponent}
If
\(\log\omega=(\log x)^{\theta+o(1)}\), then the sourced cumulative
bound proves a positive power-of-log saving for the exceptional
mass only when \(\theta+c_0>1\); the available exponent is then
\(\theta+c_0-1\).  At or below that threshold the source bound by
itself gives no decay on the terminal window.  If several endpoint
dilates or CRT pullbacks are required, the corresponding union must
be bounded before an exceptional-window exponent is entered.
\end{remark}

\begin{definition}[Pointwise interface]
\label{def:pointwise}
For an actual component family \(\mathcal F_X\), the pointwise
interface with error \(\varepsilon_X\) is
\[
 \sup_{f\in\mathcal F_X}
 \sup_{T\in I_f}
 \frac{|S_f(T)|}{V_f(T)}
 \le\varepsilon_X,
\]
after zero-mass prefixes have been discharged.  Every mask, weight,
phase, origin and prefix belongs to the theorem statement.
\end{definition}

\begin{definition}[Selector domination]
\label{def:selector}
Let \(\nu_X\) be the probability measure with which the literal
physical synthesis samples the scale variable.  It has selector
constant \(K_X\ge1\) if
\begin{equation}
 \nu_X(E)\le K_Xm_X(E)
\label{eq:domination}
\end{equation}
for every Borel \(E\subseteq J_X\).
\end{definition}

\begin{theorem}[Carleson selector return]
\label{thm:selector}
Suppose \eqref{eq:almost-scale} holds for a nonnegative measurable
function \(|C|\), and \cref{def:selector} holds.  Then
\begin{equation}
\boxed{
 \int_{J_X}|C(t)|\,d\nu_X(t)
 \le K_X\varepsilon_X.}
\label{eq:return}
\end{equation}
\end{theorem}

\begin{proof}
The measure inequality \eqref{eq:domination} is equivalent to
\(\nu_X\ll m_X\) with Radon--Nikodym derivative at most \(K_X\)
almost everywhere.  Integrate \(|C|\) against that derivative and
use \eqref{eq:almost-scale}.
\end{proof}

\begin{proposition}[Atomic selector stop]
\label{prop:atomic}
If \(\nu_X\) has an atom of positive mass, then
\eqref{eq:domination} fails for every finite \(K_X\).
\end{proposition}

\begin{proof}
For an atom \(t_0\), one has
\(\nu_X(\{t_0\})>0\), while
\(m_X(\{t_0\})=0\).
\end{proof}

Therefore a literal discrete prefix maximum cannot be called a
Carleson-distributed scale family.  One must instead prove
\cref{def:pointwise}, or prove a smoothing identity that replaces
each atom by an absolutely continuous window and bounds the
resulting stability error.

\begin{remark}[Family-level use]
\label{rem:family-selector}
\Cref{thm:selector} controls the \(\nu_X\)-average of one
nonnegative error; it does not control a supremum.  For a component
family, one needs either a uniform arithmetic error and uniform
selector constant component by component, or one joint domination
statement on component--scale space.  The total absolute component
mass is then charged exactly once as the census loss below.
\end{remark}

\section{Power transport}

Assume now that every input is expressed at amplitude scale.  Let
\[
 \varepsilon_X=X^{-\sigma_{\rm aff}+o(1)}
\]
be either a pointwise arithmetic rate or an almost-scale rate.
For the selector route put
\[
 K_X=X^{\ell_{\rm sel}+o(1)};
\]
for the pointwise route set \(\ell_{\rm sel}=0\).  Let
\(\ell_{\rm cen}\) be the theorem-backed weighted component census,
\(\ell_{\rm BV}\) the cost of returning the literal physical weight
and phase from the arithmetic coefficient, and let the squarefree
tail save \(X^{-\eta_{\rm tail}+o(1)}\).
All loss exponents in this paragraph are nonnegative and every
input is already normalized at amplitude scale.

\begin{theorem}[Typed raw exponent]
\label{thm:raw}
Under the preceding proved premises, the returned raw amplitude
saving is
\begin{equation}
\boxed{
 \sigma_{\rm raw}
 =
 \min\left\{
 \eta_{\rm tail},
 \sigma_{\rm aff}
 -\ell_{\rm sel}
 -\ell_{\rm cen}
 -\ell_{\rm BV}
 \right\}.}
\label{eq:raw}
\end{equation}
It is a positive power only if both entries of the minimum are
positive.
\end{theorem}

\begin{proof}
The tail is bounded separately.  On the retained part,
\cref{thm:selector} costs \(X^{\ell_{\rm sel}+o(1)}\), while the
weighted census and BV return cost
\(X^{\ell_{\rm cen}+\ell_{\rm BV}+o(1)}\).  Multiplying these bounds
subtracts their exponents from \(\sigma_{\rm aff}\).  Adding the
tail selects the weaker saving.
\end{proof}

\begin{corollary}[Endpoint gate]
\label{cor:endpoint}
If the complete physical occurrence loss of TPC-131 is
\(\Lambda_{\rm phys}\), then the selected route proves the final
\(o(X)\) conclusion only under the strict inequality
\begin{equation}
 \sigma_{\rm raw}>\Lambda_{\rm phys}.
\label{eq:endpoint}
\end{equation}
Equality is a stop.  If the local MVP1 raw target is fixed to
\(1/400\), every upstream loss must be included before that target
is recorded.  The exact remaining endpoint margin is
\[
 \sigma_{\rm final}
 :=
 \sigma_{\rm raw}-\Lambda_{\rm phys}.
\]
\end{corollary}

\begin{proof}
Insert \eqref{eq:raw} into the typed synthesis theorem of TPC-131
\citep{WangTPC131}.
\end{proof}

\section{The logarithmic corridor and its separate ledger}

TPC-137 gives a qualitative fixed-data logarithmic theorem
\citep{WangTPC137}.  TPC-139 shows that only a non-effective slow
growing data envelope follows from that theorem automatically, but
also records Tao--Ter\"av\"ainen's explicit small-polylog affine
estimate outside a small logarithmic-density exceptional set
\citep{WangTPC139,TaoTeravainen2026}.
After the literal component and scale normalization has been
matched, boundedness turns their good-scale estimate plus the
exceptional-set density bound into an interface of the form
\eqref{eq:almost-scale} with a power-of-log error.  That matching is
one of the eligibility checks, not an automatic property of the
physical archive.
If
\[
 \varepsilon_X=(\log\log X)^{-c}
\]
or \((\log X)^{-c}\), or merely \(o(1)\), then the largest fixed
\(X\)-power exponent certified by that statement is
\[
 \sigma_{\rm aff}=0.
\]
Indeed, for every fixed \(\sigma>0\),
\((\log\log X)^{-c}\) is eventually larger than \(X^{-\sigma}\).
Thus a logarithmic rate is real cancellation but cannot be entered
as \(1/400\).

Even a qualitative conclusion after selector transport requires
the explicit product \(K_X\varepsilon_X\to0\).  The labels
\(K_X=X^{o(1)}\) and \(\varepsilon_X=o(1)\), without rates, do not
imply that product tends to zero.

There is nevertheless a useful logarithmic ledger.  Separate the
nonexceptional correlation rate from the exceptional-window mass.
Suppose
\[
 \varepsilon_X^{\rm corr}
 =(\log X)^{-\kappa_{\rm corr}+o(1)},\qquad
 m_X(\mathcal E_{\rm tot})
 =(\log X)^{-\kappa_{\rm exc}+o(1)},
\]
where \(\mathcal E_{\rm tot}\) is the union of every required
pulled-back exceptional set, and put
\begin{equation}
 \kappa_{\rm aff}
 =
 \min\{\kappa_{\rm corr},\kappa_{\rm exc}\}.
\label{eq:effective-log-aff}
\end{equation}
Indeed, boundedness and the split into
\(\mathcal E_{\rm tot}\) and its complement give an arithmetic
scale-average error of order
\(\varepsilon_X^{\rm corr}+m_X(\mathcal E_{\rm tot})\).
Also suppose
\[
 K_X=(\log X)^{\kappa_{\rm sel}+o(1)}.
\]
Write \(\kappa_{\rm cen},\kappa_{\rm BV}\) for the
power-of-log census and return losses, while the tail saves
\((\log X)^{-\kappa_{\rm tail}+o(1)}\).  The same proof as
\cref{thm:raw} gives
\begin{equation}
\boxed{
 \kappa_{\rm raw}
 =
 \min\left\{
 \kappa_{\rm tail},
 \kappa_{\rm aff}-\kappa_{\rm sel}
 -\kappa_{\rm cen}-\kappa_{\rm BV}
 \right\}.}
\label{eq:log-ledger}
\end{equation}
If \(\kappa_{\rm raw}>0\), this closes a qualitative
almost-scale/selector return.  It still records
\(\sigma_{\rm raw}=0\) in the \(X\)-power ledger and therefore
cannot pay any fixed positive \(\Lambda_{\rm phys}\).
\Cref{prop:global-to-window,rem:window-exponent} show that
\(\kappa_{\rm exc}\) cannot simply be copied from the source's
global-density exponent for an arbitrary terminal window.

\section{Certificate schema and current verdict}

A route record must contain at least
\[
\begin{array}{ll}
\texttt{carrier}:&\text{actual fixed-two carrier hash},\\
\texttt{arithmetic}:&\sigma_{\rm aff}\text{ and theorem source},\\
\texttt{return\_mode}:&\texttt{pointwise}\text{ or }\texttt{selector},\\
\texttt{selector}:&K_X\text{ and domination source if used},\\
\texttt{losses}:&\eta_{\rm tail},\ell_{\rm cen},\ell_{\rm BV},\\
\texttt{endpoint}:&\Lambda_{\rm phys}\text{ and strict comparison},\\
\texttt{evidence}:&\text{proved, conditional, finite, or open}.
\end{array}
\]
The deterministic script checks this ledger with exact rational
arithmetic.  It includes a hypothetical strict pass solely to test
the formula, an equality case that must stop, and the current
zero-power case.  The hypothetical record is not evidence about an
actual coefficient.

\begin{center}
\small
\begin{tabularx}{\textwidth}{@{}p{0.48\textwidth}X@{}}
\toprule
Statement & Level and status\\
\midrule
Selector implication \eqref{eq:return}
& Exact measure theorem (\(\Lzero\)).\\
Atomic selector obstruction
& Exact firewall (\(\Lzero\)).\\
Raw exponent formula \eqref{eq:raw}
& Conditional typed synthesis (\(\Lone\)).\\
Logarithmic exponent formula \eqref{eq:log-ledger}
& Conditional qualitative return (\(\Lone\)); zero \(X\)-power.\\
Actual growing determinant-two power theorem
& Open.\\
Actual pointwise prefixes or selector domination
& Open / not supplied.\\
Positive fixed-\(h_0\) L2, H3, or \(1/400\)
& Not proved.\\
\bottomrule
\end{tabularx}
\end{center}

The present arithmetic subgraph therefore emits no
\textsc{arithmetic frontier} certificate.  Without a complete
actual archive it remains \textsc{not testable}; even after that
archive arrives, the first arithmetic missing item is a growing
determinant-two rate together with either a pointwise all-prefix
theorem or a valid selector crosswalk.

\begingroup
\footnotesize
\raggedright
\bibliographystyle{plainnat}
\bibliography{references}
\endgroup

\end{document}
